\documentclass[12pt]{article}
\usepackage[utf8]{inputenc}
\usepackage[T1]{fontenc}
\usepackage{hyperref}
\title{Stieng, Budeh Language Resources for K-12 Curriculum Development}
\author{Compiled by Claude Sonnet 5.5 (claude-sonnet-5-5) for Andr\'as Kornai}
\date{}
\begin{document}
\maketitle

\section{Introduction}
Budeh Stieng (ISO 639-3 \texttt{stt}, Glottocode bude1236) is an Austroasiatic language of Vietnam. Web search found little published material for this variety. Nothing found is a school curriculum resource. Be careful not to confuse it with Stieng Bulo (\texttt{sti}), a different variety.

\section{Primary Educational Resources}
No primers, textbooks or school programmes in Budeh were found in search results.

\section{Linguistic and Reference Materials}
\begin{itemize}
\item H.\ Azémar's French--Stieng dictionary (about 2,500 words collected in 1865, published in Saigon in 1887), listed in Glottolog: \url{https://glottolog.org/resource/reference/id/547146}
\item A related Azémar entry in Glottolog: \url{https://glottolog.org/resource/reference/id/19156}
\end{itemize}
No modern Budeh dictionary or grammar was found.

\section{Religious and Community Resources}
\begin{itemize}
\item Joshua Project profile for Budeh Stieng. It reports Bible portions and a New Testament (2009 or later) but no complete Bible. Search results suggest the New Testament may be a fairly literal rendering of the Vietnamese text: \url{https://joshuaproject.net/languages/stt}
\end{itemize}

\section{Digital and Audio Resources}
\begin{itemize}
\item Global Recordings Network, Budeh audio (Bible lessons and evangelistic messages): \url{https://globalrecordings.net/en/language/stt}
\item Global Recordings Network, Stieng, Budeh language page: \url{https://globalrecordings.net/ja/language/3169}
\end{itemize}

\section{Curriculum Development Context}
Little exists for Budeh. A curriculum would have to start from new orthography work, primers and community-produced materials. The audio recordings above are mainly religious and would need adaptation. The resources found are thin and none is educational.

\end{document}